\documentclass[11pt,a4paper]{article}

\usepackage[utf8]{inputenc}

\usepackage[T1]{fontenc}

\usepackage[polish]{babel}

\usepackage{lmodern}

\usepackage{geometry}

\geometry{margin=27mm}

\usepackage[hidelinks]{hyperref}

\title{Odporność ekologiczna pomiaru recovery między urządzeniami i kontekstami}

\author{Mateusz Pomianek}

\date{Szkic naukowy, wersja 1.0.0 | 28 listopada 2026}

\begin{document}

\maketitle

\begin{abstract}

Proponowana miara główna to różnica między kontekstami w medianie bezwzględnych rozbieżności czasów wznowienia.

\end{abstract}

\section{Wprowadzenie}

Analiza zgodności metod pomiaru wymaga jawnego oszacowania różnic między pomiarami (Bland \& Altman, 1986). Pomiar w różnych kontekstach wymaga rozdzielenia wpływu urządzenia od zmian badanego zachowania (Shiffman et al., 2008).

\section{Cel i pytania badawcze}

Hipotezy robocze wymagają zatwierdzenia przed rozpoczęciem właściwej weryfikacji.

\paragraph{H1.} Bezwzględna różnica czasu wznowienia wyznaczanego na dwóch urządzeniach będzie większa podczas sesji mobilnej z rozproszeniami niż w sesji o kontrolowanej pozycji i oświetleniu.

\section{Materiał i metody}

Pierwszorzędowy wskaźnik: różnica między kontekstami w medianie bezwzględnych rozbieżności czasów wznowienia.

W obu kontekstach urządzenia mają rejestrować ten sam epizod równocześnie, ze wspólną referencją czasową. Kontrast kontekstów dotyczy mediany bezwzględnych różnic, z bootstrapem po uczestnikach. Oddzielnie należy oszacować obciążenie kierunkowe i granice zgodności oraz odsetek epizodów, w których jedno urządzenie nie dostarczyło pomiaru.

Przed badaniem z udziałem ludzi należy zatwierdzić wymagania etyczne, zasady zgody i ochronę danych.

\section{Plan analizy}

Zamrożenie estymandu, reguł oceny jakości i obsługi braków nastąpi przed dostępem do danych wynikowych. Raport obejmie oszacowanie kontrastu głównego z przedziałem niepewności i pełne mianowniki. Jeśli przewidziano hipotezy drugorzędowe, przed analizą należy ustalić sposób kontroli wielokrotności.

\section{Wyniki}

W niniejszym pliku nie umieszczono wyników empirycznych, statystyk ani nieprzeprowadzonych analiz. Sekcja wymaga uzupełnienia po uzyskaniu i weryfikacji materiału dowodowego.

\section{Dyskusja i ograniczenia}

Ograniczenia obejmują możliwe błędy pomiaru, dobór próby, alternatywne mechanizmy i ograniczoną możliwość uogólniania. Korelacja nie dowodzi przyczynowości; poprawność techniczna nie jest dowodem znaczenia klinicznego.

\section{Wnioski}

Hipotezy i sposób pomiaru stanowią podstawę oceny wykonalności oraz przygotowania właściwego protokołu.

\section*{Bibliografia}

{\setlength{\parindent}{0pt}
\setlength{\parskip}{5pt}
\hangindent=1.27cm
Bland, J. M., \& Altman, D. G. (1986). Statistical methods for assessing agreement between two methods of clinical measurement. \textit{The Lancet, 327}(8476), 307--310. \url{https://doi.org/10.1016/S0140-6736(86)90837-8}\par
\hangindent=1.27cm
Shiffman, S., Stone, A. A., \& Hufford, M. R. (2008). Ecological momentary assessment. \textit{Annual Review of Clinical Psychology, 4}, 1--32. \url{https://doi.org/10.1146/annurev.clinpsy.3.022806.091415}\par
}

\end{document}
